% Generated from locale/tr/content/model-theory/basics/isomorphism.tex; do not hand-edit.


\olfileid[tr]{mod}{bas}{iso}
\olsection{İzomorf Yapılar}

Birinci dereceden !!{structure}s iki biçimde birbirine benzeyebilir.
Birincisi, aynı !!{sentence}s ifadelerini doğru kılmalarıdır. Böyle
!!{structure}s ifadelerine \emph{elementer eşdeğer} deriz. Ancak yapılar
çok farklı oldukları halde aynı !!{sentence}s ifadelerini doğru
kılabilir; örneğin biri !!{enumerable} iken öteki olmayabilir. Bunun
nedeni, dilin yeterince zengin olmaması ya da Löwenheim--Skolem teoremi
gibi birinci dereceden mantığın temel sınırlamaları yüzünden
!!a{structure} ifadesinin birçok özelliğinin birinci dereceden dillerde
ifade edilememesidir. Dolayısıyla !!{structure}s ifadelerinin
benzeyebileceği daha katı bir anlam da, onları oluşturan nesneler
dışında yapısal özellikleri bakımından farklılık göstermemeleri; yani
özünde aynı olmalarıdır. Bunu kesinleştirmenin bir yolu
\emph{izomorfizma} kavramıdır.

\begin{defn}
\ollabel{defn:elem-equiv}
Aynı !!{language}~$\Lang L$ için iki !!{structure}s~$\Struct M$ ve
$\Struct M'$ verilsin. $\Lang L$ dilinin her !!{sentence}~$!A$ ifadesi
için $\Sat{M}{!A}$ ancak ve ancak $\Sat{M'}{!A}$ ise $\Struct M$
yapısının $\Struct M'$ yapısına \emph{elementer eşdeğer} olduğunu
söyler ve $\Struct M\equiv\Struct M'$ yazarız.
\end{defn}

\begin{defn}
\ollabel{defn:isomorphism}
Aynı !!{language}~$\Lang L$ için iki !!{structure}s~$\Struct M$ ve
$\Struct M'$ verilsin. Aşağıdaki koşulları sağlayan bir
$h\colon\Domain M\to\Domain{M'}$ fonksiyonu varsa ve ancak varsa
$\Struct M$ yapısının $\Struct M'$ yapısına \emph{izomorf} olduğunu
söyler ve $\Struct M\simeq\Struct M'$ yazarız:
\begin{enumerate}
\item $h$ !!{injective}: $h(x)=h(y)$ ise $x=y$;
\item $h$ !!{surjective}: Her $y\in\Domain{M'}$ için $h(x)=y$ olacak
  bir $x\in\Domain M$ vardır;
\item \ollabel{defn:iso-const}her !!{constant}~$c$ için
  $h(\Assign{c}{M})=\Assign{c}{M'}$;
\item \ollabel{defn:iso-pred}her $n$ yerli !!{predicate}~$P$ için
  \[
  \tuple{a_1, \dots, a_n}\in \Assign{P}{M} \quad\text{ancak ve ancak}\quad
  \tuple{h(a_1), \dots, h(a_n)} \in \Assign{P}{M'};
  \]
\item \ollabel{defn:iso-func}her $n$ yerli !!{function}~$f$ için
  \[
  h(\Assign{f}{M}(a_1, \dots, a_n)) =
  \Assign{f}{M'}(h(a_1), \dots, h(a_n)).
  \]
\end{enumerate}
\end{defn}

\begin{thm}
\ollabel{thm:isom}
$\Struct M\iso\Struct M'$ ise $\Struct M\elemequiv\Struct{M'}$.
\end{thm}

\begin{proof}
$h$, $\Struct M$ yapısından $\Struct M'$ yapısına bir izomorfizma
olsun. Her atama~$s$ için $h\circ s$, $h$ ile $s$'nin bileşkesidir;
yani $\Struct M'$ içindeki $(h\circ s)(x)=h(s(x))$ atamasıdır. $t$ ve
$!A$ üzerinde tümevarımla şu daha güçlü savlar ispatlanabilir:
\begin{enumerate}
  \item[a.] $h(\Value{t}{M}[s]) = \Value{t}{M'}[h\circ s]$.
  \item[b.] $\Sat{M}{!A}[s]$ ancak ve ancak
    $\Sat{M'}{!A}[h\circ s]$.
\end{enumerate}
Birinci sav $t$'nin karmaşıklığı üzerinde tümevarımla ispatlanır.
\begin{enumerate}
\item $t\ident c$ ise $\Value{c}{M}[s]=\Assign{c}{M}$ ve
  $\Value{c}{M'}[h\circ s]=\Assign{c}{M'}$. Dolayısıyla
  $h(\Value{t}{M}[s])=h(\Assign{c}{M})=\Assign{c}{M'}$
  (\olref{defn:isomorphism} tanımındaki \olref{defn:iso-const}
  gereği) $=\Value{t}{M'}[h\circ s]$.
\item $t\ident x$ ise $\Value{x}{M}[s]=s(x)$ ve
  $\Value{x}{M'}[h\circ s]=h(s(x))$. Dolayısıyla
  $h(\Value{x}{M}[s])=h(s(x))=\Value{x}{M'}[h\circ s]$.
\item $t\ident f(t_1,\dots,t_n)$ ise
  \begin{align*}
    \Value{t}{M}[s] & = \Assign{f}{M}(\Value{t_1}{M}[s], \dots, \Value{t_n}{M}[s])
      \quad\text{ve}\\
  \Value{t}{M'}[h \circ s] & = \Assign{f}{M'}(\Value{t_1}{M'}[h \circ
    s], \dots, \Value{t_n}{M'}[h \circ s]).
  \end{align*}
  Tümevarım varsayımına göre her $i$ için
  $h(\Value{t_i}{M}[s])=\Value{t_i}{M'}[h\circ s]$. Öyleyse
  \begin{align}
    h(\Value{t}{M}[s])
    & = h(\Assign{f}{M}(\Value{t_1}{M}[s], \dots, \Value{t_n}{M}[s])) \notag\\
    & = \Assign{f}{M'}(h(\Value{t_1}{M}[s]), \dots,
    h(\Value{t_n}{M}[s])) \ollabel{iso-1}\\
    & = \Assign{f}{M'}(\Value{t_1}{M'}[h \circ s], \dots,
    \Value{t_n}{M'}[h \circ s]) \ollabel{iso-2}\\
    & = \Value{t}{M'}[h\circ s]. \notag
  \end{align}
  Burada \olref{iso-1}, \olref{defn:isomorphism} tanımındaki
  \olref{defn:iso-func} koşulundan; \olref{iso-2} ise tümevarım
  varsayımından çıkar.
\end{enumerate}
(b) kısmı alıştırma olarak bırakılmıştır.

$!A$ bir tümceyse $s$ ve $h\circ s$ atamalarının önemi yoktur;
$\Sat{M}{!A}$ ancak ve ancak $\Sat{M'}{!A}$ elde ederiz.
\end{proof}

\begin{prob}
\olref[mod][bas][iso]{thm:isom} teoreminin (b) kısmını ayrıntılı olarak
ispatlayın. \olref[mod][bas][iso]{defn:isomorphism} tanımında
izomorfizmaları niteleyen beş özelliğin her birinin nerede
kullanıldığını belirtin.
\end{prob}

\begin{defn}
Bir $\Struct M$ yapısının \emph{otomorfizması}, $\Struct M$ yapısından
kendisine bir izomorfizmadır.
\end{defn}

\begin{prob}
Herhangi bir $\Struct M$ yapısı için $X$, $\Struct M$ içinde
tanımlanabilir bir alt küme ve $h$, $\Struct M$ yapısının bir
otomorfizmasıysa $X=\Setabs{h(x)}{x\in X}$ olduğunu (yani $X$'in $h$
altında değişmez kaldığını) gösterin.
\end{prob}

